\documentclass[11pt]{article}
\usepackage{amsmath,amssymb}
\title{Reconnecting a single-file formation across an information break}
\date{24 November 2026}
\begin{document}
\maketitle

\paragraph{Known boundary.}
Use the ordered corridor, fixed anchors, speed bound, and clearance
model of \emph{Clearance and connectivity after a corridor crash}.
The disabled body separates the surviving robots into two sides.
Each side is assumed to know its own initial positions. Both sides
know the fixed obstacle position, the requirements, and when their
motion budget starts. No positions or plans are passed across the
break. Time for acquiring the assumptions and coordinating within
each side is not included in the motion budget.

Apply the interval scan from the outer anchor toward the body on each side,
reversing coordinates on the right. If either side has no admissible plan,
neither does the whole formation. Otherwise let \(x^*\) be the largest
attainable coordinate of the left inner robot and \(y^*\) the smallest on
the right. For all attainable pairs \((x,y)\), \(y^*-x^*\le y-x\).
The cross-body lower clearance is already enforced by each side
(\(y^*-x^*\ge2g\)). Thus the independently reconstructed extreme plans
reconnect by the deadline exactly when the whole configuration can.
Line interpolation preserves clearance and side membership. The argument
needs attained extrema and separate side clearance, not interval
attainable sets. Shared resources or a nonmonotone cross-gap constraint
can break the inference. Knowing the policy guarantees success is not
the same as having operational evidence of successful execution.

\paragraph{An unobserved stop changes the answer.}
Take anchors \(0,10\), movable survivors at \(2,8\), speed one,
clearance \(1\), and radio range \(7/2\). The disabled body stops at \(f=4\)
or \(f=6\). The right robot knows \(f\); the left robot cannot distinguish
the worlds from its initial view. With \(f\) known, terminal pairs
\((3,13/2)\) and \((7/2,7)\), respectively, are feasible.
If the left receives no distinguishing information by the commitment
cut, its terminal position must be at most \(3\) for safety in the first
world. In the second, the right must stay at or above \(7\), forcing a
gap of at least \(4\). No single left prescription serves both, despite
physical feasibility in each. A timely report of which of the two
positions applies would suffice; this statement does not grant that
report for free.

\paragraph{Using actual contact to create the missing observation.}
Assume instead that until a guaranteed timed receipt or its timeout,
the only readout available to the left is the passage of time.
The right follows a known policy: move as fast as allowed toward
its nearest still-safe terminal point, \(13/2\) or \(7\), and offer a
receipt at first radio contact. Receipt delivery is reliable by
\(d\ge0\) time units after contact; either its delivery before the
deadline cut is observed or its guaranteed deadline has passed.
When delivery and timeout coincide, delivery is read first. No contact
is available outside range.

The exact robust recovery time is \(2+d\).
Before observations differ, the common left motion cannot exceed \(3\).
The right cannot be left of \(8-t\) at time \(t\); therefore contact in
the \(f=4\) world cannot occur before \(3/2\). In the \(f=6\) world
the right cannot go below \(7\), leaving the gap out of range while
the left stays at most \(3\). For a worst allowed receipt delay
\(d\), no common-history prefix can separate the worlds before
\(3/2+d\). The left still needs at least \(1/2\) more unit of motion in
the \(f=6\) world. For attainment, move the left robot to \(3\) by time
one, then have it wait until \(3/2+d\). The right moves to its
specified terminal point. When \(f=4\), the receipt has arrived and
the pair is connected. Otherwise the absence after the certified cut
distinguishes \(f=6\), and the left moves to \(7/2\).
Permanent receipt erasure, added occlusion, an untrusted right policy,
or unbounded delay removes the timeout inference. Clearance
does not depend on a mere claim that a receipt was sent.

\paragraph{Verification and limits.}
The split replay checks 60,720 cases against whole-configuration
exhaustive search. The active replay checks 135 half-tick deadline and
delay pairs against exact search over the left observation histories
for the specified right policy. Lean checks the endpoint join and the
rational obstruction and timing inequalities. The source correspondence
and trajectory construction are the written arguments above. The
shared clock and timely local right estimate, trusted contact rule,
single-file geometry, and simple dynamics are assumptions.

A matched split optimizer can choose the same extreme points; a matched
contingent planner using the same physical observation models can
choose the same timeout policy. The timed case isolates how physical
motion creates usable evidence, rather than claiming improvement over
those baselines. It supplies a source-scoped R1/R6 result and an R3a
reference. It does not settle any of those tracks in full.
\end{document}
